\documentclass[12pt]{article}
\usepackage[utf8]{inputenc}
\usepackage[T1]{fontenc}
\usepackage[margin=1in]{geometry}
\usepackage{setspace}
\usepackage{xcolor}
\usepackage[colorlinks=true, urlcolor=blue!60!black, linkcolor=blue!60!black]{hyperref}
\setlength{\parindent}{0em}
\setlength{\parskip}{0.45em}
\renewcommand{\baselinestretch}{1.0}
\begin{document}
\begin{center}
{\Large \textit{Jackson v. Metropolitan Edison Co.}, 419 U.S. 345 (1974)\par}
\vspace{0.3em}
{\normalsize Case Brief --- One-Page Edition\par}
\end{center}

\textbf{Memory Jogger.} A private utility's power shutoff is not state action (\href{https://supreme.justia.com/cases/federal/us/419/345/}{419 U.S.\ 345}).

\textbf{Facts.} Metropolitan Edison, a privately owned, heavily regulated Pennsylvania utility, cut off Catherine Jackson's electricity without notice or a hearing after a billing and meter-tampering dispute, acting under a rule in its PUC-filed tariff that the commission never reviewed. She claimed a state-law entitlement to continuous service.

\textbf{Procedural History.} Her \S\ 1983 due process suit was dismissed; the Third Circuit affirmed.

\textbf{Votes.} Affirmed, 6--3 (Rehnquist for the Court; Douglas, Brennan, and Marshall dissenting).

\textbf{Holding.} State action requires a ``sufficiently close nexus'' between the State and the specific challenged act. The public-function doctrine reaches only powers ``traditionally exclusively reserved to the State,'' not utility service.

\textbf{Reasoning.} Regulation, monopoly status, and essential service were insufficient one by one, and passive acceptance of the tariff did not put the State's weight behind the shutoff, unlike \textit{Pollak}. Douglas would weigh the factors in the aggregate; Marshall warned the rule would shield the only electric company in town even if it refused service by race.

\textbf{Examples.} \emph{Same-Side:} a Medicaid nursing home's transfer (\textit{Blum}). \emph{Opposite-Side:} a utility's use of eminent domain. \emph{Fence-Sitter:} a charter school expulsion.

\textbf{Critique.} Critics fault the seriatim analysis and the treatment of commission silence as neutral; \textit{Memphis Light v.\ Craft} required due process for a municipal utility, so protection turns on ownership alone.

\textbf{Key Quotation.} ``But the inquiry must be whether there is a sufficiently close nexus between the State and the challenged action of the regulated entity so that the action of the latter may be fairly treated as that of the State itself.'' --- Justice Rehnquist, 419 U.S.\ at 351.

\end{document}
